\section*{Розділ III. Повноваження СР СМ}
\addcontentsline{toc}{section}{Розділ III. Повноваження СР СМ}

\subsection*{3.1. Представницькі повноваження}
\addcontentsline{toc}{subsection}{3.1. Представницькі повноваження}
    3.1.1. СР СМ представляє та захищає права та законні інтереси студентів Університету у побутовій сфері перед адміністрацією Університету та її структурними підрозділами, відповідальними за студентське містечко та гуртожитки.

    3.1.2. СР СМ виступає від імені студентської спільноти Університету у взаємовідносинах з комунальними службами, організаціями та іншими структурами з питань побуту, безпеки та інфраструктури студмістечка.

\subsection*{3.2. Виконавчі повноваження}
\addcontentsline{toc}{subsection}{3.2. Виконавчі повноваження}
    3.2.1. СР СМ організовує виконання рішень Конференції студентів Університету (КСУ) у побутовій сфері, забезпечуючи розробку необхідних заходів, координацію виконавців та контроль за реалізацією.

    3.2.2. СР СМ реалізує основні завдання студентського самоврядування у побутовій сфері, визначені Положенням про ОСС та цим Положенням, у межах своєї компетенції.

    3.2.3. СР СМ звітує перед КСУ про виконання покладених на неї завдань та рішень КСУ.

\subsection*{3.3. Організаційні повноваження}
\addcontentsline{toc}{subsection}{3.3. Організаційні повноваження}
    3.3.1. СР СМ організовує та проводить заходи культурного, спортивного, соціального та іншого характеру на території студмістечка, спрямовані на покращення умов проживання та розвиток студентського середовища.

    3.3.2. СР СМ координує роботу своїх структурних підрозділів (комітетів, робочих груп тощо).

    3.3.3. СР СМ забезпечує інформаційну підтримку мешканців гуртожитків та студентів Університету щодо побутових питань, діяльності ОСС у побутовій сфері та важливих подій у студмістечку.

\subsection*{3.4. Повноваження щодо погодження рішень}
\addcontentsline{toc}{subsection}{3.4. Повноваження щодо погодження рішень}
    3.4.1. СР СМ реалізує право органів студентського самоврядування на участь в управлінні Університетом шляхом погодження рішень адміністрації Університету, що стосуються побутової сфери, у випадках та порядку, передбачених Законом України "Про вищу освіту", Статутом Університету та Положенням про ОСС.

\subsection*{3.5. Повноваження щодо СРГ}
\addcontentsline{toc}{subsection}{3.5. Повноваження щодо СРГ}
    3.5.1. СР СМ здійснює координацію діяльності Студентських рад гуртожитків (СРГ), забезпечуючи єдність підходів до реалізації завдань студентського самоврядування у побутовій сфері.

    3.5.2. СР СМ надає методичну та організаційну допомогу СРГ.

    3.5.3. СР СМ може розробляти та надавати СРГ рекомендації щодо організації їхньої роботи та проведення заходів.

    3.5.4. СР СМ здійснює контроль за діяльністю СРГ в межах повноважень, визначених Положенням про ОСС та/або делегованих КСУ.

    3.5.5. СР СМ реалізує повноваження щодо тимчасового управління справами студентського самоврядування гуртожитку у випадках та порядку, визначених п. 6.8.2 цього Положення.

\subsection*{3.6. Контрольні повноваження}
\addcontentsline{toc}{subsection}{3.6. Контрольні повноваження}
    3.6.1. СР СМ здійснює внутрішній контроль за діяльністю своїх структурних підрозділів та посадових осіб.

    3.6.2. СР СМ взаємодіє зі Студентською уповноваженою делегацією (СУД) з питань контролю за дотриманням нормативних документів ОСС та цільовим використанням ресурсів у межах своєї компетенції.

\subsection*{3.7. Обов'язковість рішень СР СМ}
\addcontentsline{toc}{subsection}{3.7. Обов'язковість рішень СР СМ}
    3.7.1. Рішення СР СМ, прийняті колегіально в межах її повноважень відповідно до цього Положення та Положення про ОСС, є обов'язковими для виконання всіма членами СР СМ, її структурними підрозділами (комітетами, робочими групами тощо) та посадовими особами СР СМ.

    3.7.2. Невиконання або неналежне виконання рішень СР СМ є підставою для відповідальності відповідно до цього Положення та інших внутрішніх документів ОСС.
